\section{What Version Information Actually Buys}
\label{sec:consequences}

The Gaussian core lets us separate three kinds of version information that are
often conflated in regression testing: localization of where a change may occur,
statistical coupling of old and new executions, and numerical slack in the old
fairness certificate.  The first two can change query complexity directly.  The
third, by itself, does not improve the worst-case rate when patch effects are
unrestricted.

\subsection{A linear separation from patch localization}

Consider a fairness graph that is a matching of $M$ counterfactual edges and
assume unit paired-delta variance at every vertex.  Compare two testers with the
same paired execution oracle.  A \emph{patch-aware} tester receives a sound
impact set identifying the unique edge on which behavior may have changed.  A
\emph{location-unaware} tester is only promised that at most one edge changed,
but is not told which edge.

\begin{corollary}[Linear patch-locality separation]
\label{cor:locality-separation}
Fix $\delta\in(0,1/2)$ and $\gamma>0$.  On the $M$-edge matching family above:
\begin{enumerate}[label=(\roman*)]
\item a patch-aware tester uses at most
\[
\frac{32\log(2/\delta)}{\gamma^2}+2
\]
paired-oracle calls; while
\item every adaptive location-unaware $\delta$-correct tester with finite expected
stopping time uses, on the all-neutral null instance,
\[
\E_0[T]\ge
\frac{8M\,\kl(1-\delta,\delta)}{\gamma^2}.
\]
\end{enumerate}
Consequently, for every fixed $\delta$, the worst-case separation is
$\Omega(M)$.
\end{corollary}

\begin{proof}
With a revealed sound impact edge, the relevant graph $H$ consists of one edge,
for which \Cref{prop:families} gives $\Psi(H)=4$.  Substitution into
\Cref{thm:upper} gives the first claim.

Without the edge location, correctness must hold against an alternative on every
one of the $M$ disjoint edges.  The relevant testing graph is therefore the full
matching, for which \Cref{prop:families} gives $\Psi=4M$.  Applying
\Cref{thm:lower} yields the second claim.  The ratio is linear in $M$ for fixed
confidence.
\end{proof}

The comparison does not give the patch-aware algorithm a stronger execution
oracle: its only additional information is the sound impact location.  Thus the
separation quantifies, in a clean family, the statistical value of software
change-impact information.

\subsection{Old certificate slack is not enough in the minimax sense}

Suppose an old edge $e=(u,v)$ is certified with gap
$g_{0,e}=\mu_{0,u}-\mu_{0,v}$ satisfying $|g_{0,e}|\le\tau_e$.  It is tempting to
expect large headroom $\tau_e-|g_{0,e}|$ to reduce the worst-case cost of testing
the new version.  This intuition is false without a restriction on the patch.

\begin{proposition}[No minimax gain from old slack alone]
\label{prop:slack-no-free-lunch}
Consider a single fairness edge with Gaussian paired-delta observations of known
variances $\nu_u^2,\nu_v^2>0$.  The old gap $g_0$ and threshold $\tau$ may be known
exactly, with arbitrary slack $|g_0|\le\tau$.  A tester must distinguish
\[
\mathsf H_0:\ |g_0+\theta_u-\theta_v|\le\tau
\]
from
\[
\mathsf H_1:\ |g_0+\theta_u-\theta_v|\ge\tau+\gamma.
\]
For every adaptive $\delta$-correct tester with finite expected stopping time,
there is a boundary null instance for which
\[
\E[T]\ge
\frac{2\kl(1-\delta,\delta)}{\gamma^2}
\min_{q_u,q_v>0}
\left\{q_u+q_v:
\frac{\nu_u^2}{q_u}+\frac{\nu_v^2}{q_v}\le1
\right\}.
\]
In particular, the lower bound is independent of the old certificate slack.
\end{proposition}

\begin{proof}
Choose a sign $s\in\{-1,+1\}$ and a null patch effect whose edge difference is
\[
\theta_u^{(0)}-\theta_v^{(0)}=s\tau-g_0,
\]
so the new program lies exactly on the fairness boundary.  Construct an
alternative by increasing the signed edge difference by $\gamma$:
\[
\theta_u^{(1)}-\theta_v^{(1)}=s(\tau+\gamma)-g_0.
\]
The two hypotheses therefore differ by exactly $\gamma$ in the endpoint mean
difference, regardless of $g_0$.  Splitting this $\gamma$ shift adversarially
between the two endpoints and applying the same adaptive change-of-measure
calculation as in \Cref{thm:lower} gives the stated bound.
\end{proof}

\Cref{prop:slack-no-free-lunch} does not say that historical certificates are
useless.  They can yield instance-dependent savings when the observed patch stays
well inside the boundary, and they become powerful under structural assumptions
that limit patch magnitude or direction.  The proposition says that such savings
cannot be claimed as a uniform minimax consequence of old slack alone.


\subsection{A two-axis information hierarchy}

The preceding Gaussian separation isolates the value of knowing \emph{where} a
patch may have changed behavior, while \Cref{thm:bernoulli-separation} isolates
the value of knowing \emph{how to compare} the two versions through replayable
randomness.  The two effects can be combined on a matching family.

Let \(G\) consist of \(M\) disjoint counterfactual edges
\(e_i=(u_i,v_i)\).  Fix \(p\in(0,1)\) and \(d\in(0,1-p)\).  The old version and
the null new version have Bernoulli parameter \(p\) at every vertex.  Alternative
\(i\) changes only \(u_i\) from \(p\) to \(p+d\).  Thus exactly one fairness edge
may regress, but its location is unknown.

\begin{corollary}[Unknown-location replay complexity]
\label{cor:bernoulli-location}
Under common-uniform paired replay, a tester with no patch-location information
can be \(\delta\)-correct using
\[
M\left\lceil
\frac{\log\delta}{\log(1-d)}
\right\rceil
\]
paired-oracle calls.  Conversely, every adaptive \(\delta\)-correct tester with
finite expected stopping time satisfies on the all-neutral null
\[
\E_0[T]
\ge
\frac{M\,\kl(1-\delta,\delta)}
{(1-p)\log\!\left(\frac{1-p}{1-p-d}\right)}.
\]
Consequently, for fixed interior \(p\),
\[
Q_{\rm pair}^{\rm unknown}(M,d,\delta)
=
\Theta\!\left(
\frac{M\log(1/\delta)}{d}
\right)
\]
up to universal constant factors for, say, \(\delta\le 1/4\).
If a sound impact analysis identifies the affected edge in advance, the factor
\(M\) disappears by \Cref{thm:bernoulli-separation}.
\end{corollary}

\begin{proof}
For the upper bound, query each candidate vertex \(u_i\) for
\[
n=\left\lceil\frac{\log\delta}{\log(1-d)}\right\rceil
\]
paired replay rounds and reject if any old/new discordance is observed.  Under the
null, common-uniform replay makes every paired outcome identical, so the type-I
error is exactly zero even though \(M\) locations are inspected.  Under
alternative \(i\), only the samples at \(u_i\) are informative and each is
discordant with probability \(d\); hence the miss probability is
\((1-d)^n\le\delta\).

For the lower bound, fix alternative \(i\).  As in
\Cref{thm:bernoulli-separation}, a paired query to \(u_i\) contributes
\[
K_{\rm pair}
=
(1-p)\log\!\left(\frac{1-p}{1-p-d}\right)
\]
null-to-alternative KL, while every query to every other vertex has identical law
under the two hypotheses and contributes zero.  Adaptive transcript
change-of-measure therefore implies
\[
\E_0[N_{u_i}]K_{\rm pair}
\ge
\kl(1-\delta,\delta).
\]
The \(M\) candidate vertices are distinct.  Summing these inequalities over
\(i=1,\ldots,M\) yields the stated lower bound.
\end{proof}

\Cref{cor:bernoulli-location} should not be read as a claim that searching for one
abnormal stochastic source is new; active sequential anomaly search is an
established information-theoretic problem \cite{CohenZhao2015}.  Its role here is
to expose a software-regression decomposition: sound change-impact information
removes the multiplicative location-search cost, while replayable old/new
randomness changes the small-patch dependence from quadratic to linear.
